\begin{frame}{Discussão}
  \begin{itemize}
    \item[$\blacktriangleright$] Análise por contexto de crédito
      \begin{itemize}
        \item \textbf{Crédito Pessoal}
          \begin{itemize}
            \item Técnicas de re-amostragem, especialmente o RUS, apresentaram desempenho superior.
            \item Random Forest e XGBoost foram os modelos mais eficazes
          \end{itemize}
        \item \textbf{Crédito Corporativo}
          \begin{itemize}
            \item SMOTE se destacou como a técnica mais eficaz
            \item Desempenho geral foi baixo, devido ao pequeno número de amostras e à alta dimensionalidade dos dados.
          \end{itemize}
      \end{itemize}
  \end{itemize}
\end{frame}

\begin{frame}{Discussão}
  \begin{itemize}
    \item[$\blacktriangleright$] Desempenho Custo-sensível vs Re-amostragem
      \begin{itemize}
        \item Em geral, as estratégias de re-amostragem (RUS, SMOTE e SMOTE-Tomek) apresentaram desempenho superior em comparação com o CS-SVM.
        \item O CS-SVM mostrou-se competitivo em alguns casos, mas não conseguiu superar consistentemente as técnicas de re-amostragem.
      \end{itemize}
  \end{itemize}
\end{frame}

\begin{frame}{Discussão}
  \begin{itemize}
    \item[$\blacktriangleright$] \textit{Trade-off} entre Precisão e Sensibilidade
      \begin{itemize}
        \item As técnicas de re-amostragem, especialmente o RUS, tendem a melhorar a sensibilidade às classes minoritárias, mas isso pode ocorrer às custas da precisão.
        \item O CS-SVM apresentou um equilíbrio mais favorável entre precisão e sensibilidade em alguns casos, mas não foi suficiente para superar as técnicas de re-amostragem.
      \end{itemize}
  \end{itemize}
\end{frame}
